\section{The limiting nonunit margin profile}\label{tks:sec:margins}
Choose uniformly a symmetric sorted-positive-margin matrix counted by $a_n$. For $j\geq2$, let $M_j$ be the number of its margins equal to $j$. A profile $m=(m_j)_{j\geq2}$ has finite support and is feasible when
\begin{equation}\label{tks:eq:marginweight}
W(m)=\sum_{j\geq2}jm_j\leq n.
\end{equation}
It then determines the remaining $n-W(m)$ unit margins. This profile refers to distinct row positions in the matrix model, not to a quotient by simultaneous permutations.

Put
\begin{equation}\label{tks:eq:pilaw}
w(m)=\prod_{j\geq2}(j!)^{-m_j},\qquad
\pi(m)=\frac{w(m)}{C},\qquad
T_3(m)=\sum_{j\geq3}2\binom j3m_j.
\end{equation}
The law $\pi$ is the product of independent geometric distributions on $\{0,1,2,\ldots\}$,
\begin{equation}\label{tks:eq:geommargin}
\PP_\pi(M_j=m)=(1-1/j!)(1/j!)^m.
\end{equation}
Since $\sum_{j\geq2}\PP_\pi(M_j>0)=\sum_{j\geq2}1/j!<\infty$, it is supported on finite profiles. Its normalization is exactly $C$, and
\begin{equation}\label{tks:eq:ET3}
\E_\pi T_3=\sum_{j\geq3}\frac{2\binom j3}{j!-1}=H_3.
\end{equation}

Let $P_n$ be the actual profile law, extended by zero to infeasible profiles, and put $q_n=I_{n-3}/I_n$ for $n\geq3$.
\begin{theorem}\label{tks:thm:marginlaw}
On the countable space of finite nonunit margin profiles,
\begin{equation}\label{tks:eq:marginL1}
\left\|\frac{\dd P_n}{\dd\pi}-1-q_n(T_3-H_3)\right\|_{L^1(\pi)}
 =O(n^{-2}).
\end{equation}
In particular,
\begin{equation}\label{tks:eq:marginTV}
\TV(P_n,\pi)=\frac{2H_3}{C}n^{-3/2}+O(n^{-2}).
\end{equation}
Thus the entire nonunit margin profile converges jointly in total variation to the independent geometric law.
\end{theorem}
\begin{proof}
For a feasible profile, apply the Young-subgroup sum in \eqref{tks:eq:schwob} to its single margin partition. The identity permutation contributes exactly $I_nw(m)$. No support-$1$ or support-$2$ element contributes. A support-$3$ element is a single $3$-cycle in a block of size $j$, with $2\binom j3$ choices per such block. Its square-root count is $I_{n-3}$. Therefore the number of matrices of profile $m$ is exactly
\begin{equation}\label{tks:eq:fixedprofile}
I_nw(m)+I_{n-3}w(m)T_3(m)+R_{n,m},\qquad R_{n,m}\geq0,
\end{equation}
where the remainder contains supports at least $4$. The proof of Theorem~\ref{tks:thm:main} gives
\begin{equation}\label{tks:eq:marginremainder}
\sum_{W(m)\leq n}R_{n,m}=O(I_n n^{-2}).
\end{equation}

For every $1<R<\sqrt2$, the exponential moment
\[
\E_\pi R^W
 =\prod_{j\geq2}\frac{1-1/j!}{1-R^j/j!}
\]
is finite. The $T_3$-weighted moment is also finite, either by differentiating the geometric factors and summing a factorially decaying series, or by bounding $T_3$ by a polynomial in $W$ and choosing a slightly larger admissible $R$. Consequently both
$\sum_{W>n}\pi(m)$ and $\sum_{W>n}\pi(m)T_3(m)$ are $O(R^{-n})$.
Summing \eqref{tks:eq:fixedprofile} and using \eqref{tks:eq:marginremainder} now gives
\[
\frac{a_n}{CI_n}=1+q_nH_3+O(n^{-2}).
\]
Divide \eqref{tks:eq:fixedprofile} by $a_n\pi(m)$, sum absolute errors against $\pi$, and include the exponentially small infeasible tail. The remainder has $L^1$ mass $O(n^{-2})$ by positivity. Since $q_n=O(n^{-3/2})$, quadratic normalization terms are $O(n^{-3})$. This proves \eqref{tks:eq:marginL1}.

Total variation is half the $L^1$ density difference, so
\[
\TV(P_n,\pi)=\frac{q_n}{2}\E_\pi|T_3-H_3|+O(n^{-2}).
\]
The variable $T_3$ is nonnegative and every positive value is at least $2$. The certificate in Section~\ref{tks:sec:constants} proves $0<H_3<2$. Since its mean is $H_3$,
\[
\E_\pi|T_3-H_3|=2H_3\PP_\pi(T_3=0),\qquad
\PP_\pi(T_3=0)=\prod_{j\geq3}(1-1/j!)=\frac2C.
\]
Finally \eqref{tks:eq:shiftratio} gives $q_n=n^{-3/2}+O(n^{-2})$, proving \eqref{tks:eq:marginTV}.
\end{proof}

\subsection{The matrix dimension deficit}
If $K_n$ is the matrix dimension, then
\[
n-K_n=\sum_{j\geq2}(j-1)M_j.
\]
Total variation contracts under a measurable map. Thus the law of this deficit converges, with error at most that in \eqref{tks:eq:marginTV}, to the finite random variable
$D=\sum_{j\geq2}(j-1)M_j$ under \eqref{tks:eq:geommargin}. Its probability generating function is
\begin{equation}\label{tks:eq:deficitpgf}
\E z^D=\prod_{j\geq2}\frac{1-1/j!}{1-z^{j-1}/j!},\qquad |z|\leq1.
\end{equation}
For example, the limiting probability that every margin equals $1$ is $1/C$, and that no margin exceeds $2$ is $2/C$. These follow from the joint profile theorem. Distributional and total-variation convergence do not by themselves assert convergence of unbounded moments of the actual finite-$n$ deficit; no such additional assertion is needed here.
